\section{Analysis Models}\label{app:models}

\subsection{Request State Diagram}

\figref{fig:states} shows the states of a Request and every transition between them; Section~\ref{feat:lifecycle} gives the guards in full. Three properties of the diagram carry the design: a Hold exists only while the Request is Pending, so no Hold outlives the Objection Window; Charges and Earnings post only on the transitions into Allocated; and Waiting is the single state from which a Request recovers, whether it lost a line to a refusal, a lapse, or a Bump.

\begin{landscape}
  \centering
  \includegraphics[width=\linewidth,height=0.86\textheight,keepaspectratio]{request-states.pdf}
  \captionof{figure}{Request state-transition diagram. Submitted is a transient evaluation state left in the same transaction as it is entered. Withdraw is drawn once, from Draft; a Manager may also withdraw from Pending and Waiting. Allocated, Withdrawn, and Refused are final, except that Allocated returns to Waiting when a Required line is Bumped.}
  \label{fig:states}
\end{landscape}

A typical Request with all three classes runs Draft $\to$ Submitted $\to$ Pending (room and GPUs on Hold, engineers awaiting approval) $\to$ Allocated (last approval arrives; Holds commit; Charges and Earnings post). The two failure paths are: a refusal, after which the Manager picks a replacement and the Request returns to Pending if the Objection Window is still running; and a lapse, after which the Capacity Planner decides the outstanding approvals and the Request is Allocated only when every Required line can be satisfied at once, with no Hold in between.

\subsection{Use Case Diagram}

\figref{fig:usecases} shows the twelve use cases of Appendix~\ref{app:usecases} and the actors that start them. The Clock actor stands for the scheduled jobs that fire when an Objection Window lapses or a Hold expires; it has no user interface.

\begin{figure}[htbp]
  \centering
  \includegraphics[width=\textwidth,height=0.85\textheight,keepaspectratio]{use-cases.pdf}
  \caption{Use case diagram. All human actors are Members; the roles are additive.}
  \label{fig:usecases}
\end{figure}

\subsection{Context Diagram and Data Model}

The context diagram is \figref{fig:context} in Section~\ref{sec:overall}; the data model is \figref{fig:datamodel} in Appendix~\ref{app:data}.
